\documentclass[11pt]{article}
% Shared preamble for the round 1 response letters.
% RED-D-26-00078, submitted as "Housing First or Treatment First?",
% resubmitted as "Housing First or Treatment First? Ask the Kids, Not the Parents"
%
% House style, decided once here:
%   * responses are in RED, referee text is in the black italic block above them
%   * NO HYPHENS anywhere in the prose
%   * every response ends with a pointer to where in the new paper it lives

\usepackage[margin=1.05in]{geometry}
\usepackage{amsmath,amssymb}
\usepackage{xcolor}
\usepackage[official]{eurosym}
\usepackage{booktabs}
\usepackage{longtable}
\usepackage{array}
\usepackage[hidelinks]{hyperref}

% enumitem is not in this TeX installation, so first level list spacing is set
% directly instead
\makeatletter
\renewcommand{\@listI}{%
  \leftmargin=1.7em \labelwidth=1.2em
  \topsep=0.4em \parsep=0pt \itemsep=0.3em \partopsep=0pt}
\makeatother

\definecolor{respred}{RGB}{190,25,25}

\setlength{\parskip}{0.35em}
\setlength{\parindent}{0pt}

\newcolumntype{L}[1]{>{\raggedright\arraybackslash}p{#1}}
\newcommand{\pname}[1]{\texttt{#1}}

% ---- the two building blocks -------------------------------------------------
% \cmt{label}{short title} opens a comment, then the verbatim or paraphrased
% referee text goes in the referee environment, then the response.

\newcommand{\cmt}[2]{%
  \par\medskip
  \noindent\textbf{#1\quad #2}\par}

\newenvironment{referee}
  {\begin{list}{}{\leftmargin=1.6em \rightmargin=1.2em \topsep=0.4em}\item[]\small\itshape}
  {\end{list}}

\newenvironment{response}
  {\par\begingroup\color{respred}\medskip}
  {\par\endgroup\medskip}

% a flag for statements in the reports that are factually wrong, so they are
% easy to find again and easy to handle carefully
\newcommand{\correction}{\textbf{[Point of correction.]}\ }
\newcommand{\concede}{\textbf{[Conceded.]}\ }
\newcommand{\newresult}{\textbf{[New result.]}\ }
\newcommand{\open}{\textbf{[Still open.]}\ }

% numbers that are not in yet, so they are impossible to miss on a read through
\newcommand{\tbd}[1]{\textbf{[TO FILL: #1]}}

% Marks what has been done since the first draft of these letters. The letters
% were written when most of this was a plan; this flag marks the places where it
% is now built and running.
\newcommand{\thisround}{\textbf{[Built since the first draft of this letter.]}\ }

\newcommand{\lettertitle}[2]{%
  \begin{center}
  {\Large\bfseries Housing First or Treatment First?}\\[0.3em]
  {\normalsize Manuscript RED-D-26-00078}\\[0.9em]
  {\large\bfseries #1}\\[0.3em]
  {\normalsize #2}
  \end{center}
  \vspace{0.6em}}


\begin{document}

\lettertitle{Response to Referee 3}{Major and minor concerns, responses in red}

Thank you for the report, and there is no need to apologise for the timing. Your
question at M2(c), where the case for intervention comes from if drug use is
rational, is now the question the paper is organised around.


{\color{respred}
\textbf{One note before the comments.} The paper has a new title: it was submitted as
\emph{Housing First or Treatment First?} and is now \emph{Housing First or Treatment First? Ask the Kids, Not the Parents}. The reason is the substance of the revision rather than a relabelling, and it
will be clearest at M2(c), which asks where the externalities come from.

\textbf{What changed that matters for your comments.} The paper now asks whether
housing assistance, in the programme as it exists and therefore rationed, should
require sobriety, and it begins by asking whether moral hazard alone justifies the
condition. It does not. With every effect of a parent's drug use on its child
removed, a newborn is worse off under Treatment First, by 1.09 percent of lifetime
consumption, even when the savings of the screen (Treatment First's drug test, which
revokes the voucher of a user it catches) buy more vouchers for other households.
The answer turns instead on the children, through three externalities that a parent
counts only through an imperfect weight on its child. A parent's drug spending and
lost earnings are consumption and housing its child goes without (the budget
externality); a child whose parent had a substance use disorder (SUD) while it lived
at home starts adult life with lower productivity (the earnings externality); and a
child lives wherever its parent lives (the childhood homelessness externality). In
the long run, with the drug price clearing its market and the voucher probability
balancing the budget, Treatment First cuts drug use by 21 percent and homelessness by
11 percent, and a newborn gains 2.57 percent of lifetime consumption, three quarters of
it an externality its parent does not count. The parents gain too, by 0.69 percent of
their own lifetime consumption. The transition is now solved for every cohort alive at
the reform, with the drug price and the budget clearing every period (Sections~6.4
and~6.5). Along the main transition, homelessness and childhood homelessness never
rise, and children born at the reform gain 1.45 percent.

Several of the changes answer your comments directly. Preferences are the same in
every state, so reducing the disorder no longer raises the utility of consumption by
construction (M2(a)). A new two period model (Section~3) shows that it is the
uninsurable mortality risk of a SUD that puts drug use at the bottom of the income
distribution, and the quantitative model's distribution of use by income, which is
not targeted, is now compared with the data (M1(c), M2(b)). The drug price clears a
supply curve, and the demand elasticities are checks rather than targets (M2(d)). The loss of
consortium is gone (m2). Two things you asked about are no longer in the paper, and
I say why where they come up: labour supply is now inelastic (m1), and the
experiments no longer move from rationed to universal vouchers (m3).
}

%======================================================================
\section*{Major concern 1: connection to empirical evidence and stylized facts}

\cmt{M1(a)}{Present stylized facts prominently}
\begin{referee}
``The nature of the model, with its large number of parameters, makes it
relatively opaque. I believe it would be helpful for the paper to include a few
stylized facts more prominently, notably on the empirically observed correlations
between drug use, voucher receipt, and homelessness.''
\end{referee}
\begin{response}
You are right that the facts were hard to see behind the model. I have not added a
separate empirical section. Instead, each fact now sits where it is used, next to the
model's counterpart, so the reader sees both at once.

\begin{itemize}
\item \textbf{Drug use and income.} Section~3 opens with the two facts the stylized
model is built to explain. A user in the paper is an adult with a drug use disorder of
any severity, mild included, and fewer than one in ten adults are users: 8.90 percent
of adults aged 26 and over in the 2024 National Survey on Drug Use and Health
(NSDUH). And the average user's income is 0.74 of the population average. Figure~2
now shows the whole relationship (M1(c)).
\item \textbf{Drug use and homelessness.} A SUD, meaning a moderate or severe
disorder, is far more common among the homeless than among all adults: 21.7 percent of
the homeless have one (Gutwinski et al., 2021), against 3.68 percent of adults aged 26
and over in the 2024 NSDUH. The same holds for serious mental illness (SMI): 25
percent of the homeless, against 4.82 percent of adults in the 2024 NSDUH. The shares
among the homeless are calibration targets (Table~3, Blocks~2 and~3). Table~4 adds the age
profile of the homeless, how many of them work and what they earn (Block~1).
\item \textbf{Voucher receipt.} 26.6 percent of income eligible households hold a
voucher in the Picture of Subsidized Households of the Department of Housing and Urban
Development (HUD) for 2023, the target for the probability of drawing one (Table~3,
Block~2). Table~4, Block~3, adds what holders pay in rent and how many live mainly on
wages.
\end{itemize}

\open One correlation you name is not in the paper: the paper does not report drug
use among voucher holders, in the data or in the model. On vouchers and homelessness,
it cites the trials in which housing assistance reduces homelessness (Section~2), but
it does not compare the model with them.

The Associate Editor asked (A1) for a stylized model whose predictions align with a
few stylized facts. Section~3 does this for drug use and income, which decides whose
children bear the costs of drug use.

The facts are in the opening of Section~3, in Figure~2 (Section~5.2, Block~2), and in
Tables~3 and~4.
\end{response}

\cmt{M1(b)}{Survival and unemployment probabilities are not transparent}
\begin{referee}
``Survival probabilities (especially their dependence on homelessness and drug
use) and unemployment probabilities are an essential input into the model's
calibration and welfare implications. It would be good to present more prominently
some empirical evidence on their correlation with homelessness and drug use, which
are currently borrowed from other papers but not always set transparently (for
example for unemployment transition probabilities).''
\end{referee}
\begin{response}
Both are now set out in one place in the model, Section~4.8, which takes the death
risk, the unemployment risk and the child's productivity penalty in turn. Every
multiplier has its source in Table~3 or Table~A1, with the model's value beside the
data wherever the parameter is calibrated.

\textbf{Survival.} Baseline survival by age is the 2010 life table for the total
population of the United States (Arias, 2014), for ages 26 to 85 (Appendix~A,
Block~10). Three multipliers raise the probability of dying, equation~(20):
\begin{itemize}
\item a SUD, by $\kappa_{\phi d}=4.00$. It is calibrated to the income of users
relative to the average, because it is the parameter that decides how far down the
income distribution use sits (Section~3.3). The value lies within the mortality rate
ratios of a SUD in the literature: 3.91 in the Danish registers (Plana-Ripoll et al.,
2019) and 6.6 among Norwegians in specialist care for a SUD (Heiberg et al., 2018).
\item homelessness, by $\kappa_{\phi H}=1.60$ at every age, the mortality ratio of
homeless adults against poor housed adults in Meyer et al.\ (2023).
\item SMI, by $\kappa_{\phi m}=2.50$, the middle of the standardised mortality ratios
for SMI in Olfson et al.\ (2015).
\end{itemize}
Two untargeted checks bear on them (Table~4). Excess deaths due to a SUD are 14.5
percent of deaths at ages 25 to 44 in the model, against 20.0 percent from drug
poisonings in the WONDER database of the Centers for Disease Control and Prevention
(Block~2). And for households with SMI whom Treatment First's revocation puts on the
street, losing the voucher raises three year mortality by about 3.6 percentage points
in the model, against the 4.6 points by which enrolling in Housing First lowers it
among homeless veterans with mental illness in Costantini (2025) (Block~7).

\textbf{Unemployment.} Equation~(21) sets out the transitions. The baseline separation
rate is calibrated to the aggregate unemployment rate, 4.9 percent over 2014 to 2024
in the Current Population Survey (5.6 percent in the model), and is tilted towards low
productivity households with the unemployment rates by education of the Bureau of
Labor Statistics. The job finding rate is 0.55, from Krueger et al.\ (2014) (Appendix~A, Block~9). A SUD and SMI
do not raise separations. Each lowers the chance of finding work, by a divisor
calibrated to the share of the homeless with that condition, 21.7 percent for a SUD
and 25 percent for SMI (Gutwinski et al., 2021), since a household out of work longer
is likelier to lose its home (Table~3, Blocks~2 and~3). Homelessness makes separations
20 percent more frequent and job finding 20 percent slower, from Desmond and
Gershenson (2016).

The unemployment rate given a SUD is no longer a target. It is measured in the 2024
NSDUH for adults aged 26 and over, the ages the model covers, at 16.2 percent, and the
model gives 14.5 percent without being fitted to it (Table~4, Block~2). The paper
also reports where the model misses: 21.4 percent of the homeless are employed in the
model, against 46 percent in Meyer et al.\ (2024), and their income is 6.5 percent of
the average, against 19.1 (Table~4, Block~1).

All of this is in Section~4.8, in Tables~3, 4 and~A1, and in Appendix~A, Blocks~9
and~10.
\end{response}

\cmt{M1(c)}{The income to drug use relationship is not shown}
\begin{referee}
``\dots the relationship between income and drug use (severe or mild) is key to
the calibration of $b$ and the income-dependent penalties in the model \dots but
its empirical counterpart does not feature very prominently in the paper. It would
be helpful to get a sense of this joint distribution in the data.''
\end{referee}
\begin{response}
Done, and it is now a figure. Figure~2 (Section~5.2, Block~2) plots the share of
households that use drugs at each level of income relative to the mean, in the model
and in the 2024 NSDUH. The calibration targets only the average income of users, so
the rest of the figure is a check. As in the data, use is highest at the bottom and
lowest at the top: 18.3 percent of households with less than 0.22 of mean income use
drugs, against 13.7 percent in the data, and 4.2 percent of those with more than 1.62
times the mean, against 4.8 percent. Between 0.43 and 1.62 times the mean, use in the
model is flatter and lower than in the data, 5.9 to 6.3 percent against 7.2 to 10.2.
Both differences lower the average income of users in the model, which is why it falls
slightly short of its target: users earn 62 percent of the average in the model,
against 74 percent in the data (Table~3).

The paper says why the age and the income of users both matter for the experiments
(Block~2). Both policies reach a child through its parent, so the age of users decides
how many children their use touches. And the vouchers are means tested, so the lower
users sit in the income distribution, the more of them, and of their children, the
screen reaches.

On severity, a user is a household whose use amounts to a disorder of any severity,
mild included, and a household with a SUD is one whose disorder is moderate or severe.
Use that never reaches a disorder is not modelled. Figure~2 is drawn for users in this
sense. \open The paper does not show the gradient separately for moderate or severe
disorders.

What the gradient disciplines has changed. In the previous draft it was key to $b$
and to penalties that depended on income, as you say. There are no income dependent
penalties now: a SUD changes the same hazards by the same factor at every income. The
average income of users pins one parameter, the mortality multiplier of a SUD
$\kappa_{\phi d}$, and $b$ is set to the participation rate, 8.90 percent. The reason
is the result of the stylized model (Section~3.3, Figure~1). Drugs are a luxury good bought at a price, so
without any consequence of use the average user would earn 3.51 times the population
average. The mortality risk of a SUD alone takes this to 0.64 at its calibrated
strength, because dying cannot be insured against and what it destroys grows with
income. The unemployment risk alone leaves it at 3.51, because a household can save
against it; with saving switched off, the unemployment risk alone takes it to 0.75.
The income of users in the data is therefore informative about the mortality channel
in particular, and that is the parameter it is used for.

This is in Section~3.3 and Figure~1, in the $\kappa_{\phi d}$ item of Section~5.1
(Block~1), and in Section~5.2 (Block~2, Figure~2).
\end{response}

%======================================================================
\section*{Major concern 2: modelling of the drug addiction process}

\cmt{M2(a)}{Reducing the disorder mechanically raises the utility of consumption}
\begin{referee}
``\dots reducing addiction lowers $\theta_d$ \dots and thus increases directly
consumption utility from housing and non-housing, thus implying mechanically that
policies reducing addictions and policies raising housing consumption are strong
welfare complements even absent incentive effects. I am not sure there is an
alternative model that would make more sense, but I would like the author to
acknowledge and potentially provide more sensitivity to this modelling choice.''
\end{referee}
\begin{response}
\concede You are right, and Referee~2 raises the same point (C5). The weight on the
drug no longer changes with the disorder. Flow utility, equation~(12) in Section~4.2,
is the same function of consumption, housing and the drug in every state: $\theta$ and
$b$ are the same with and without a SUD, with and without SMI, and housed or homeless,
and SMI does not enter utility at all (Section~4.7). The stylized model has the same
property (equations~(2) and~(3)), and no preference parameter varies with income
either.

Reducing the disorder therefore no longer raises the utility of consumption and
housing by itself. A SUD makes a household worse off only through its consequences in
the model: a higher risk of dying, a lower chance of finding work, a higher chance of
SMI, the earnings penalty its child inherits, which the parent counts at its weight on
the child, and, under Treatment First, the risk of losing its voucher unless it is in
treatment and does not use. The one term of flow utility that depends on the disorder
is the cost of a year in treatment, which only a household with a SUD can incur, and
only if it chooses treatment.

The complementarity between housing and recovery is now a transition probability, not
a taste. Recovery is less likely while homeless, by the multiplier $\mu_H$, which
targets the effect of housing on abstinence in Milby et al.\ (2005), 1.47 in the data
against 1.45 in the model (Table~3, Block~4). The same trial randomised housing
conditional on abstinence against housing without the condition, and the model gives
1.40 against its ratio of 1.26 without targeting it (Table~4, Block~8).

One interaction remains, and I should say so. The drug is in the same bundle as
consumption and housing, equation~(11), so with risk aversion of two, how much a
household uses still changes what consumption and housing are worth to it at the
margin. This comes from the quantity used, in the same way for every household, and not
from the disorder. With log utility the bundle separates and the interaction
disappears; the qualitative results of the stylized model, on who uses drugs, hold with
log utility, recalibrated to the same targets (a footnote in Section~3.1). \open The
paper does not recompute the welfare results of the quantitative model under another
functional form.

The preferences are in Section~4.2, equations~(11) and~(12).
\end{response}

\cmt{M2(b)}{What patterns does the Stone Geary form actually generate?}
\begin{referee}
``I find the use of the Stone-Geary functional form to generate a large mass of
non-users, combined with income-dependent penalties \dots very clever. However, I
am unclear on what are the actual patterns of drug consumption this generates, and
where heterogeneity in $v$ (conditional on addiction status $d$) stems from. It
would be very helpful to present model-derived stylized facts on drug use \dots by
income, employment status, housing status, assets \dots and explain how they depend
on model parameters (notably $b$, the Stone-Geary parameter, and $\aleph$, the value
of life).''
\end{referee}
\begin{response}
The new Section~3 answers the first half of this, and Section~5.2 the second.

\textbf{How the functional form works.} The marginal utility of the drug is
proportional to $\theta/(b+v)$ (Section~3.1). With $b>0$ the marginal utility of the
first unit is finite, so a household with little to spend buys none. $b$ therefore
mostly decides whether a household uses at all, and $\theta$ how much a user uses,
that is, the share of its spending that goes on the drug. The risk of a SUD is the
same whatever the quantity (equations~(4) and~(17)), so it sits on the decision to use,
not on how much.

\textbf{Where heterogeneity in $v$ comes from, given $d$.} Among households that use,
how much they use moves with what they have to spend, through $\theta$: income, assets
and the voucher. Whether they use at all moves with what using risks, and that differs
across households with the same $d$: a household with more to lose from dying is
deterred more, as in M1(c). Across values of $d$, the difference is the risk itself. A
household without a SUD that uses risks acquiring one, while one that already has a
SUD risks nothing further by using in the benchmark: no hazard in equations~(17)
to~(21) depends on the quantity used, and the only one that depends on whether a
household uses at all, the onset of a SUD in~(17), applies to households without one.
Under Treatment First, using also risks the voucher.

\textbf{What the model generates.} The paper reports these patterns against the data:
\begin{itemize}
\item by income: Figure~2 (M1(c));
\item by age: 67 percent of users are 26 to 34 in the model, against 45 percent in
the NSDUH for 2023, and their mean age is 34.6 against 39.6, so users in the model are
younger than in the data (Table~4, Block~2);
\item by housing: 20.2 percent of the homeless have a SUD in the model, against 21.7
percent in the data, a target (Table~3, Block~2);
\item by mental health: households with SMI are 2.64 times as likely as others to use,
against 2.44 (Table~4, Block~2);
\item by employment: 14.5 percent of households with a SUD are unemployed, against
16.2 percent (Table~4, Block~2);
\item by parenthood: 5.82 percent of children live with a parent who has a SUD,
against 4.6 percent in the 2024 NSDUH (Table~4, Block~6), and 16.35 percent live with
a parent who uses drugs (Table~5).
\end{itemize}
\open The paper does not report drug use by assets or by employment status (only the
unemployment rate of households with a SUD), and it does not show how these patterns
move with $b$ or $\aleph$.

\textbf{$b$ and $\aleph$.} $b$ is set to the participation rate, 8.90 percent; as it
falls, the first unit is worth more and more households use (Section~5.1, Block~1).
$\aleph$ is calibrated, as in the previous draft, to a value of a statistical life
(VSL), now 80 times average labour income (Viscusi and Aldy, 2003); the model gives 81
(Table~3). As it rises, dying costs more and use falls. Dying destroys the whole value
of being alive, which grows with income, so the mortality risk of a SUD deters richer
households more (Section~3.3). Because $\aleph$ enters utility additively and utility
is concave, the VSL in the model rises roughly with the square of consumption, as in
Hall and Jones (2007); the paper sets this beside estimates of the income elasticity
of the VSL of 0.5 to 0.6 across people and 1.5 to 1.7 over time (a footnote in
Section~5.1). Figure~1, in the stylized model, is the closest the paper comes to the
comparative statics you ask for: it scales the mortality risk of a SUD from zero to
twice its calibrated strength, with $b$ recalibrated at every point so that only who
uses changes, not how many.

This is in Sections~3.1 and~3.3 and Figure~1, in Section~5.1 (Block~1), and in
Section~5.2 (Block~2, Figure~2 and Table~4).
\end{response}

\cmt{M2(c)}{Where do the externalities come from if drug use is rational?}
\begin{referee}
``Since drug use is `rational' and agents are forward-looking, the model should
make clearer where externalities arise that generate welfare-improving government
interventions. There are potentially many sources of fiscal externalities \dots as
well as potentially less obvious non-fiscal ones \dots It would be helpful to make
those appear analytically or numerically, and establish how optimal policies change
when some of this channels are shut down.''
\end{referee}
\begin{response}
You are right that if households are rational and forward looking, a case for
intervention needs something they do not count. The paper looks for it in two places
and finds it in the second, and the title changed because of the answer.

\textbf{Moral hazard and the budget.} Housing First insures a household against losing
its home when drug use leads to a SUD, unemployment and homelessness, so part of the
cost of using falls on the government's budget. In the model that cost runs through
the taxes a household stops paying when a SUD keeps it out of work, the unemployment
insurance (UI), food assistance and vouchers it draws, and the treatment the government
pays for (equation~(26)). These cover the lost taxes on your list and, through the
transfers and treatment it draws, the fiscal cost of a household with a SUD. They are
the case critics of Housing First usually make. Section~6.2 tests them alone. I remove
the child from welfare and from the parent's preferences and make its productivity
independent of what the parent did, so that a household's drug use can harm only
itself and the budget. Taxes are held, and the budget is balanced by the probability
of drawing a voucher, $\pi_V$. Treatment First then cuts drug use by 28 percent, and
yet a newborn loses 1.17 percent of lifetime consumption with the drug price and
$\pi_V$ held, 1.09 percent when $\pi_V$ spends the screen's savings on more vouchers,
and 0.74 percent when the drug price clears as well (Table~5). When $\pi_V$ clears, 58
percent of households gain, but only by 0.05 percent of consumption on average, while
96 percent of those with a SUD lose. The fiscal externality is real, but correcting it
is worth less than the option the screen takes from households that choose optimally.

\open Some items on your list are not in the model. The government's budget has no
separate service cost per homeless household: a shelter is free to the household, and
equation~(26) has no term for it. The loss of consortium is gone (m2). And apart from
what it does to the user's own child, a household's drug use has no direct effect on
other households. The paper does not discuss a fiscal cost of homelessness. Because
Treatment First lowers homelessness in the long run and does not raise it along the
main transition (Sections~6.2 and~6.4), I would expect such a cost to add to what the
screen saves rather than take from it, but the paper does not test this. A disamenity
from the size of the homeless population, which Referee~2 asked about, is in a footnote
to Section~2: at the value Loubi\`ere et al.\ (2020) find, it moves no long run gain in
Table~5 by more than 0.14 percentage points.

\textbf{The children.} What decides the comparison is the parent's imperfect altruism.
The parent values its child's flow utility at $\omega_c(\xi(j)-1)$ (equation~(12)), and
what its disorder does to the child's adult value at $\omega_c$, through the charge
$D(y,d)$ (equations~(24) and~(25)). The welfare criterion counts the child at its full
weight, $\xi(j)-1$ (equation~(31)), so the share $1-\omega_c$ is external. $\omega_c$
is 0.085, the weight at which removing the parent's concern for its child raises drug
use in the model by as much as the 2024 NSDUH shows: unmarried adults living with one
other adult and no child are 17.5 percent more likely to have used illicit drugs in the
past year than lone parents living with one child, reweighted to the same age, sex,
marital status and family income (Section~5.1, Block~5). Three externalities run
through that gap:
\begin{itemize}
\item \emph{the budget externality}: the child shares its parent's consumption and
dwelling but not the drug, so every unit the parent spends on the drug, and every unit
of earnings it loses while a SUD keeps it unemployed, is consumption and housing the
child goes without;
\item \emph{the earnings externality}: a child whose parent had a SUD while it lived at
home starts adult life with productivity lower by $\delta_d=0.24$ (Nilsson, 2017),
keeps it for life, and passes a lower draw to its own child (equation~(22));
\item \emph{the childhood homelessness externality}: a child lives wherever its parent
lives. This is the one critics of Treatment First expect it to make worse. In the long run
it falls, from 1.03 to 0.92 percent of children in a shelter, and it does not rise
along the main transition.
\end{itemize}

\newresult \textbf{Numerically.} Section~6.3 splits a newborn's gain exactly into four
parts (equation~(32) and Table~6). In the main experiment the newborn gains 2.57 percent
of lifetime consumption: 1.50 percent from the childhood its parent does not count
(the budget and childhood homelessness externalities), 0.46 percent from better
productivity at entry that its parent does not count (the earnings externality), 0.43
percent from its own adult life under the policy, and 0.19 percent that its parent
does count. Three quarters of the gain is external. The parents gain as well, by 0.69
percent of their own lifetime consumption.

Table~5 shuts channels down much as you propose, for the two policies rather than for
an optimal one. Its columns switch on the budget ($\pi_V$) and the drug price one at
a time, and each step is run with and without the externalities on children:

\begin{center}\small
\begin{tabular}{@{}lcccc@{}}
\toprule
Newborn's gain, percent of lifetime consumption & Both held & $\pi_V$ clears & Price clears & Both clear \\
\midrule
Without the externalities & $-1.17$ & $-1.09$ & $-0.89$ & $-0.74$ \\
With the externalities & $2.82$ & $6.45$ & $0.98$ & $2.57$ \\
\bottomrule
\end{tabular}
\end{center}

``Both held'' keeps the drug price and $\pi_V$ at their Housing First levels; without
the externalities the child appears only as an adult, so the gain is its value at 26.
The externalities turn a loss into a gain in every column. Spending the screen's
savings on vouchers is worth little without them and a great deal with them, because
the savings are larger ($\pi_V$ rises to 9.24 percent rather than 6.81) and the extra
vouchers house children. The drug price works against Treatment First: the screen
lowers demand, the price falls, and the cheaper drug draws users back, which lowers
$\pi_V$ and takes the gain from 6.45 to 2.57 percent.

The three externalities are switched on together, not one at a time. Table~6
separates the earnings externality from the two that act in childhood, the budget and
childhood homelessness externalities, which it reports together. Two single changes
are made in the text. With the earnings penalty at zero, Treatment First still raises
welfare (a footnote in Section~2). With $\omega_c$ at 0.69, the upper end of the
literature (Nishiyama, 2000), the newborn gains 3.18 percent and 23 percent of it is
external (a footnote in Section~5.1, Block~5).

I keep households rational, as the paper did before. With present bias, a sobriety
condition could be worth something to a household as a commitment device even with no
externality at all (Section~2), so the assumption works against Treatment First.

\open The paper does not solve for an optimal programme. It compares the two policies
the debate is about and shows which channel decides that comparison.

This is in Sections~4.2 and~4.10 for the parent's weight, Section~6.2 for moral hazard
alone, and Section~6.3, equation~(32) and Tables~5 and~6 for the decomposition.
\end{response}

\cmt{M2(d)}{How is the price elasticity targeted when $p_v$ is normalised to one?}
\begin{referee}
``I am somewhat unclear how the elasticity of drug use to its price is targeted,
given that the price of drugs $p_v$ is normalized to 1. The author says they vary
the price of drugs in the model, but this requires suppressing this normalization
and potentially affects the definition of other prices in the model?''
\end{referee}
\begin{response}
In the previous draft the price was exogenous and normalised to one, and the
participation elasticity was targeted by changing that price in simulations of the
benchmark (its footnote~6). The draft did not explain how the normalisation and the
change in price fit together, and it should have. The paper now separates the two.

\textbf{The price is an equilibrium price.} It clears the drug market against a supply
curve, $p_v=\nu_1 S_v^{1/\varepsilon_s}$ (equation~(27), Section~4.12). The
normalisation is confined to the scale $\nu_1$, which is set so that the price is one
in the benchmark and then held fixed across experiments. It normalises units, not
behaviour (Appendix~A, Block~13). Each experiment either holds the price at its
benchmark level or lets it clear (Table~5), and the transition lets it clear every
period.

\textbf{Moving it touches no other price.} Consumption has a price of one in the
budget constraint, equation~(16); the rent per unit of housing is held at its benchmark
value in every simulation; the interest rate is exogenous; and earnings are
productivity times an age profile. No other price in the model is defined relative to
the drug. The suppliers are outside the economy, so what they earn along the supply
curve enters neither welfare nor the budget (a footnote in Section~4.12).

\textbf{The demand elasticities are no longer targeted.} They are the response of drug
use to its price in the benchmark, and they are checks (Section~5.2, Block~5, and
Table~4). The participation elasticity is $-1.32$, against about $-0.3$ to $-0.9$ in
Saffer and Chaloupka (1999) and Grossman and Chaloupka (1998), and the intensive
elasticity among users is $-0.63$, against about $-0.3$ to $-0.8$ in Grossman and
Chaloupka (1998) and Olmstead et al.\ (2015). Participation is more elastic in the
model than in the data, and that works against Treatment First: as Treatment First
lowers the drug price, more households return to using, and the more elastic
participation is, the more of them do. The response to the screen itself, at a given
price, is a separate object, matched to the deterrence of the US military's random
drug testing (Mehay and Pacula, 1999) through the detection rate (Section~6).

\textbf{The supply elasticity.} $\varepsilon_s=1.46$, the estimate of Mace et al.\
(2020) for cannabis on Washington State administrative data (Section~5.1, Block~6). I
know of no direct estimate for an illicit drug. Cannabis is the closest market I could
find: it is still a Schedule I substance federally, so its suppliers operate under
federal prohibition, much as an illicit supplier does. It is a short run estimate,
over a month; longer windows give higher elasticities, and studies of illicit drugs
take the long run supply elasticity to be infinite (Becker et al., 2006). The lower the
elasticity, the more the price falls when Treatment First lowers demand, and the more
users it draws back (Block~6), so if the long run elasticity is higher, a short run
value works against Treatment First. At 1.46 the
price falls by 5 percent in the main experiment.

This is in Section~4.12, equation~(27); Section~5.1, Block~6; Appendix~A, Block~13;
and Section~5.2, Block~5, with Table~4.
\end{response}

%======================================================================
\section*{Minor concerns}

\cmt{m1}{The labour supply elasticity should matter more}
\begin{referee}
``I am surprised that there is not a clearer role for the labor supply elasticity
in welfare calculations, as it should matter quite substantially for the welfare
costs of balancing the budget with extra labor taxes.''
\end{referee}
\begin{response}
You are right that if a policy is financed by a higher labour income tax, the labour
supply elasticity matters for its welfare cost, and the version you read did both:
labour supply was elastic, and every experiment balanced the budget with the labour
income tax. \open The paper no longer has that margin: labour is supplied
inelastically when employed (Section~4.1). A footnote in Section~4.4 gives the reason.
I solved earlier versions of this model, the one you read among them, with endogenous
labour supply, rents, homeownership and house prices; these margins had a negligible
effect on the results, so I removed them. The Associate Editor asked in the same round
what role labour supply plays in the framework (A3), and suggested simplifying the
secondary dimensions of the model.

The financing has changed as well. No experiment raises taxes: the tax parameters are
held at their benchmark values in every experiment, and the budget is balanced by the
probability of drawing a voucher, $\pi_V$ (Section~4.11; Appendix~A, Block~11). And
Treatment First does not need financing. Users who lose their vouchers leave the
programme, and tax revenue rises, mainly because more children start adult life with
higher productivity and to a lesser extent because fewer adults are unemployed owing to
a SUD, so $\pi_V$ rises, from 6.75 to 7.80 percent in the main experiment
(Section~6.2). The deadweight cost of a higher labour tax therefore does not arise in
any experiment in the paper.

The labour market still matters, through employment rather than hours: a SUD lowers
the chance of finding work, and homelessness lowers it and raises separations
(Section~4.8).

This is in Section~4.1, the footnote in Section~4.4, Sections~4.11 and~6.2, and
Appendix~A, Block~11.
\end{response}

\cmt{m2}{Present sensitivity as graphs rather than tables}
\begin{referee}
``Some of the results of the counter-factual experiments might be more helpfully
presented in the form of graphs comparing outcomes and sensitivity to parameters,
rather than tables, for example, for the sensitivity to the loss of consortium
spillover, or to the value-of-life parameter $\aleph$.''
\end{referee}
\begin{response}
The loss of consortium is gone. Referee~1 argued (3(a) and 3(b)) that an adjustment
made after the model is solved is inconsistent with it, and that the conservative
value of a parameter without a source is zero. I agree, and the paper no longer has
one. The harm to others that decides the comparison is now inside the model, in the
parent's imperfect weight on its child (M2(c)).

The transition, where outcomes move over time, is a figure: in Figure~3, six panels
set the path against both steady states year by year, and two show welfare by year of
birth and the share of households worse off. I kept the long run (Table~5), the
decomposition (Table~6) and the alternative transitions (Table~7) as tables, because
each compares a handful of discrete experiments.

\open The paper does not plot results against $\aleph$, or against any other
parameter of the quantitative model. $\aleph$ is calibrated to a VSL of 80 times
average labour income (Viscusi and Aldy, 2003), and the paper does not vary it. The one
figure that varies a parameter continuously is Figure~1, in the stylized model, which
scales the mortality and unemployment risks of a SUD from zero to twice their
calibrated strength and shows who uses at each point. The sensitivity checks of the
quantitative model each move one thing to one alternative value, and the paper
reports them in the text, mostly in footnotes:
\begin{itemize}
\item the parent's weight on its child at 0.69 instead of 0.085, with $b$ recalibrated to
drug use: the newborn gains 3.18 percent instead of 2.57 (Section~5.1, Block~5);
\item the detection rate at 10 percent instead of 20: the newborn gains 1.65 percent
instead of 2.57 (Section~6);
\item the earnings penalty at zero: Treatment First still raises welfare (Section~2);
\item a homelessness disamenity at the value of Loubi\`ere et al.\ (2020): no long run
gain moves by more than 0.14 percentage points (Section~2);
\item households with a SUD unable to keep their voucher by entering treatment: the
newborn gains at least as much (Section~6.2);
\item each alternative transition at two detection rates, 20 and 100 percent
(Table~7);
\item a child whose parent dies before it leaves home sent into adult life at the
parent's death, rather than valued as the average child of its cohort: Treatment First
scores lower, by an amount the paper does not report, because that convention credits
a child for losing its parent early and Treatment First keeps more parents alive (a
footnote in Section~6.1).
\end{itemize}

These are in Figure~1, Figure~3, Tables~5 to~7, and the footnotes listed above.
\end{response}

\cmt{m3}{Vary the voucher receipt probability continuously}
\begin{referee}
``It would be helpful to provide a more continuous experiment that varies the
probability of voucher receipt rather than go straight for full rationing to fully
universal provision. This would help establish whether there is a Laffer curve of
voucher generosity, or whether the corner solution of universal provision is indeed
optimal, and what are the key parameters governing this.''
\end{referee}
\begin{response}
\open The paper no longer compares rationed with universal vouchers, so the sweep you
suggest is not in it. I dropped that comparison because the policy question changed.
The July 2025 executive order and the HUD funding notice that followed it, since struck
down in court (Section~1), are about whether the assistance that exists should require
treatment and sobriety, not about how much of it there is. That assistance is rationed:
26.6 percent of income eligible households hold a voucher. Every experiment therefore keeps the programme
rationed and asks whether to condition it (Section~6).

The voucher probability is not fixed, though. Unless an experiment holds it, $\pi_V$
adjusts to keep the government's net revenue at its benchmark level, so it moves as
far as the reform's savings allow (Section~4.11): from 6.75 to 7.80 percent in the main
experiment, to 9.24 percent with the drug price held, and between 6.55 and 6.90 percent
in the first decade of the transition (Section~6.4). Comparing the columns of Table~5
where $\pi_V$ is held and where it clears shows what that extra generosity is worth.
With the drug price held, spending the screen's savings on vouchers raises a newborn's
gain from 2.82 to 6.45 percent; without the externalities, the same step moves it only
from a loss of 1.17 percent to a loss of 1.09. Whether there is a Laffer curve, or
universal provision is optimal, is a question about the scale of the programme, and
the paper does not answer it.

This is in Section~4.11, Section~6.2 and Table~5, and Section~6.4.
\end{response}

\end{document}
